%--- Chapter 8 ----------------------------------------------------------------%
\chapter{Conclusion}


This dissertation has presented my contributions as part of the ATLAS collaboration. The work has focused on long-lived particle searches and on supporting the development and construction of the ITk upgrade in preparation for the HL-LHC era.



In particular, this work contributed to two complementary search strategies. Within the Disappearing track search, two $\stau$ signal models were incorporated, enabling the search to put limits on a model that is previously unexplored in ATLAS. 
This search is designed to use tracklets originating from long-lived charged particles that decay within the Inner Detector volume. The analysis reconstructs these tracklets and, for the shortest tracklets, searches for an associated soft pion originating from the decay at the end of the tracklet. 
This included searches in four channels, three layer with soft pion, three layer without soft pion, four layer with high MeT, and four layer with medium MeT. The  observed data is consistent with the Standard Model predictions to within 2$\sigma$, allowing limits to be set on the SUSY scenarios. For the $\stau$ models, exclusion limits at 320~\GeV for lifetime of 1ns in the CMSSM \stau coannihilation scenario and 300 \GeV for a lifetime of 1ns in the GMSB \stau scenario. 


In addition, a new analysis strategy combining disappearing track signatures with high ionization energy loss measurements was developed. This approach will extend the sensitivity to the heavy long-lived particle phase space. The inclusion of ionization energy loss provides an additional discriminating handle that significantly reduces fake backgrounds and enables more robust data driven background estimation techniques. The search is still under development, with expected sensitivity of 500 \GeV for 1ns lifetime in the GMSB \stau scenario. 

This thesis has contributed to the development and preparation of the ITk for HL-LHC operation. This includes work on module testing, integration procedures, and validation of detector performance under realistic operating conditions. In particular, studies of thermal cycling and mechanical stability demonstrated the robustness of pixel modules under extreme temperature variations, providing confidence in the long-term reliability of the detector. 
%This work also includes the step up of multiple test stands for the use 
These efforts support the successful production, assembly, and deployment of the ITk Inner System, which is critical for maintaining tracking performance in the high pileup environment of the HL-LHC.

Together, these contributions advance both the experimental techniques and detector capabilities required to search for physics beyond the Standard Model. As the LHC continues operation, the methods and detector developments presented in this thesis will play a role in the continued success of the ATLAS experiment and the search for new physics.

